\scanpage{21}
% lines 1–2
Calculant la loi de crochet sur $E$, on trouve

% lines 3
\[
\deleted{\unsure{X},}[\widetilde H,\widetilde X]=2\widetilde X,\ [\widetilde H,\widetilde Y]=-2\widetilde Y,\ [\widetilde X,\widetilde Y]=\widetilde H \tag*{(70)}
\]

% lines 4–8
qui est la loi bien connue du crochet de l'algèbre de Lie $\mathfrak{sl}(2)_S$ spéciale linéaire, qui est l'algèbre de Lie de $\mathrm{SL}(2)_S$. Quant à $f$, c'est l'homomorphisme d'Algèbres de Lie correspondant à l'\deleted{\unsure{hom|homom.}} \inserted{isogénie} bien connue

% lines 9
\[
\mathrm{SL}(2)_S\ \unsure{\mathbin{-}|\longrightarrow}\ \mathrm{GP}(1)_S \tag*{(71)}
\]

% lines 10–16
\deleted{qui \unsure{fa|f}} de noyau $\mu_2$ \inserted{(qui est le centre de $\mathrm{SL}(2)_S$)}. Quant à l'opération de $E'$ sur $E$, elle correspond à l'opération bien connue de $\mathrm{GP}(1)_S$ sur $\mathfrak{sl}(2)_S$, déduite par passage au quotient \deleted{\unsure{p}}de la représentation adjointe de $\mathrm{GL}(2)$ sur lui-= (qui est bien triviale sur le centre $\mu_2\ldots$). En termes des bases précédentes, on aura

% lines 17
\deleted{(71) $\rho_X=\widetilde\rho$}

% lines 18–20
\[
\left\{\begin{aligned}
\rho_E(X)=\widetilde\rho_E(\widetilde X):\quad&
\widetilde X\mapsto0,\ \widetilde Y\longmapsto\widetilde H,\ \widetilde H\longmapsto-2\widetilde X\\
\rho_E(Y)=\widetilde\rho_E(\widetilde Y):\quad&
\widetilde X\mapsto-\widetilde H,\ \widetilde Y\mapsto0,\ \widetilde H\mapsto2\widetilde Y\\
\rho_E(H)\ (=\tfrac12\widetilde\rho_{\widetilde H}\ \text{sic}):\quad&
\widetilde X\to\widetilde X,\ \widetilde Y\mapsto-\widetilde Y,\ \widetilde H\mapsto0
\end{aligned}\right. \tag*{(72)}
\]

% lines 21–22
(ici le plus commode était d'utiliser la ``formule oubliée'' (56 bis) \ldots)

% lines foot
% archive sheet number: 20
